\subsection{The Fork in the Road: Open vs.\ Closed Source in Practice}

The theoretical framework presented so far is already observable in the current AI industry landscape.

\subsubsection{Model A: The Open-Source Express---GLM and DeepSeek}

China's open-source LLM ecosystem has developed rapidly. Models such as Zhipu AI's GLM series and DeepSeek allow any enterprise to directly download model weights, deploy on proprietary servers, and fine-tune for specific business needs. This path is characterized by \textbf{disintermediation}: no third-party approval required, data sovereignty, predictable total cost of ownership, and customization freedom.

This model has particular value for non-digital-native sectors such as manufacturing, agriculture, and county-level administration. An agricultural enterprise in a county town does not need to pass a Silicon Valley private equity fund's AI audit to deploy an open-source model for supply chain optimization---the industrial-level manifestation of our claim that spectrum ascent entry remains open to all.

\subsubsection{Model B: The Closed-Source Gatekeeper---OpenAI and Private Equity AI Auditing}

In contrast, a novel business model is emerging in the closed-source ecosystem. OpenAI has reportedly partnered with private equity funds to conduct AI transformation readiness assessments of portfolio companies. Those passing the assessment gain prioritized access to OpenAI tools and additional investment. This model can be characterized as \textbf{re-intermediation}: AI capability requires external authorization; tool access is bundled with financing eligibility; the enterprise's judgment about its own AI readiness is outsourced; and data must pass through third-party servers.

This model is rational from a financial perspective, but from the perspective of cognitive democratization, it constitutes the cognitive stratification solidification mechanism we have warned against: AI capability becomes not a universal factor of production, but a privilege requiring permission.

\subsubsection{What the Contrast Teaches Us}

The coexistence of these two models gives our theoretical argument real-world urgency. The open-source path demonstrates that disintermediated AI empowerment is feasible. The closed-source path demonstrates what happens without open-source alternatives: AI capability is re-enclosed by financial intermediaries. This also provides an empirical response to the market-will-solve-accessibility counterargument: the market may reduce closed-source AI prices, but price reduction is not disintermediation. A cheaper AI that still requires a PE fund audit remains gatekeeper-controlled. Cognitive democratization requires not cheaper permission, but permissionless capability.
